\chapter{Conclusion and Future Work}
\label{chap:conclusion}

\section{Summary of the Study}

This study investigated how representative preamble-based timing
synchronization methods affect timing-estimation accuracy and end-to-end
OFDM detection. A modular MATLAB simulation framework was developed
and verified progressively through eight experiments. Baseline verification
and controlled timing displacements established the behavior of the OFDM
link before the four timing estimators were compared. Individual
estimator checks preceded the unified timing experiment (EXP07) and the
end-to-end BER experiment (EXP08).

The formal comparison used a 64-subcarrier QPSK--OFDM system with a
16-sample cyclic prefix, equal CP-inclusive training duration and energy,
and a common useful-symbol timing reference and decision interval.
Each unified experiment contained 10,000 trials for each combination of
four methods, two channel conditions and eleven SNR levels, giving
880,000 timing-estimation attempts per experiment. EXP08 additionally
evaluated three receiver configurations, giving 2,640,000 receiver
observations. The shared trial scenarios supported controlled comparisons,
while EXP07 and EXP08 used distinct random seeds. Timing MSE and
correct-estimation probabilities were the primary evaluation measures;
BER provided complementary receiver-level evidence.

\section[Research Questions]{Answers to the Research Questions}

\subsection*{RQ1: Effects of timing offset}

The baseline BER agreed closely with the theoretical coherent QPSK
reference in AWGN. The controlled-offset experiments demonstrated that
accurate FFT-window placement is essential, but that the effects of timing
offset depend on the channel and receiver processing. For a causal channel
with maximum delay $\nu$ and cyclic-prefix length $G$, the sufficient
CP-safe interval is $[\nu-G,0]$ under the adopted timing convention.
Within this interval, an integer timing displacement preserves the circular
convolution structure but produces a subcarrier-dependent phase rotation.
Outside the interval, interference or loss of the intended samples may
remain even after ideal phase correction. For example, EXP02 showed that
a five-sample early displacement was almost perfectly correctable in the
noiseless symbol-domain experiment, whereas a five-sample late displacement
retained nonzero distortion. Thus, CP protection and correct detection
are distinct properties.

\subsection*{RQ2: Comparison of the four estimators}

The unified timing results showed that metric shape, decision rule and
error distribution jointly determine statistical performance. The selected
Schmidl--Cox (S\&C) midpoint rule retained a systematic early bias associated
with the plateau. At 20 dB AWGN, its exact probability was 0.0006,
whereas its CP-safe probability was 0.9978. Its MSE was 59.4513 samples
squared, approximately 97.36\% of which was attributable to squared
bias. Low noise therefore improved safe-window placement without
eliminating the bias relative to the useful-symbol boundary.

The Park-type and Minn-type methods reduced plateau ambiguity and
achieved much higher exact probabilities at moderate and high SNR.
Nevertheless, occasional large errors could dominate their timing MSE.
The CAZAC-based known-template estimator produced the strongest
empirical timing results: all 110,000 AWGN estimates in EXP07 were
exact, and only three of its 110,000 fixed-multipath estimates were
nonexact, each by one sample. This result is conditional on the selected
waveform, prior information and bounded search interval. The comparison
concerns the implemented designs rather than every version of the
published algorithms.

\subsection*{RQ3: Sensitivity to noise and multipath}

The influence of SNR was not uniform across the estimators. In EXP07,
Park and Minn generally became more accurate as SNR increased, whereas
the S\&C midpoint remained systematically early. The specified multipath
profile was particularly challenging for peak-based estimation at low SNR.
At 0 dB in multipath, Park had a higher exact probability than Minn
(0.4856 versus 0.4653), but its MSE was much larger (322.0781 versus
96.4519 samples squared) and its CP-safe probability was lower (0.5363
versus 0.7605). The choice of performance criterion therefore changes the
relative interpretation of these two designs.

These findings describe AWGN and the particular deterministic three-tap
channel used in the study. The optional investigations of carrier-frequency
offset (CFO) and phase noise in the proposal were not undertaken.
Consequently, the answer to RQ3 is restricted to the completed noise
and multipath experiments.

\subsection*{RQ4: Relationship between timing accuracy and BER}

EXP08 demonstrated that the relationship depends critically on the receiver's
treatment of residual timing error. At 20 dB AWGN, S\&C placed 99.81\%
of its estimated windows within the CP-safe interval, yet its uncompensated
BER was approximately 0.499553. Oracle phase compensation reduced the
bit-error count from 639,428 to one out of 1,280,000 decoded bits. In
fixed multipath at the same SNR, the safe-window fraction was 99.05\%,
and the error count fell from 637,574 to 22 after compensation. In both
conditions, all remaining oracle bit errors occurred in CP-unsafe frames.
The paired frame-level analysis supported the substantial reduction without
requiring independent bit decisions within a frame.

Park, Minn and CAZAC generally approached their respective perfect-timing
BER references as SNR increased. Their references in multipath were
slightly different because the received active-frame power determined the
noise scaling separately for each training design. The appropriate
synchronization penalty is therefore interpreted relative to each method's
own reference. The oracle uses the true timing error, and its performance
quantifies the benefit of ideal phase information rather than that of a
practical phase-estimation algorithm.

\section{Contributions and Implications}

The principal contribution is a controlled statistical comparison of four
representative timing-synchronization implementations, supported by a
verified OFDM framework and saved trial records.
The separation of exact timing, CP-safe timing, MSE and estimator
validity provides a more informative evaluation than a single success
criterion. Error distributions and the bias--variance decomposition explain
why persistent small displacements and rare large errors can lead to
different conclusions even when aggregate probabilities appear favorable.
The receiver comparison and paired BER analysis extend these findings to
data detection. Zero observed errors are interpreted as finite-sample
results, rather than as proof of zero underlying error probability.

The study also highlights a trade-off between timing performance,
receiver-side information and processing cost. Direct correlation at a
single candidate requires 32 complex products for S\&C and Park, 48
for Minn, and 80 for CAZAC under the implemented parameter settings.
These counts describe core correlation workload rather than measured
latency. CAZAC's favorable accuracy is accompanied by a longer correlation
and use of the complete known template; S\&C's shorter correlation and
high safe-window fraction may be useful when effective channel estimation
or residual phase tracking is available. Park and Minn offer alternative
peak-based designs whose relative merits depend on the channel and the
chosen evaluation measure. No universally optimal estimator follows from
this comparison.

\section{Limitations and Future Work}

The conclusions are restricted to integer-sample acquisition in a
frame-present, simulated baseband link with known channel responses,
a single fixed multipath profile, and a bounded search interval. No
practical channel estimator, fractional timing recovery, noise-only
packet detection or hardware implementation was evaluated. The Park-type,
Minn-type and CAZAC-based implementations are representative variants,
not strict reproductions of every cited publication. Equal training
overhead does not ensure identical receiver-side information. These
limitations motivate the following directions for future research.

\begin{enumerate}
\item \textbf{Broader channel and impairment models.}
Extend the comparison to random fading channels, varying path strengths
and delay spreads, fractional-sample delays, CFO and phase noise.
Investigate both timing acquisition and subsequent tracking, including
situations in which the CP is insufficient for the channel delay spread.

\item \textbf{Practical receiver compensation.}
Replace oracle information with pilot-assisted or training-based effective
channel estimation and residual phase correction. Compare the resulting
BER with the existing perfect and oracle references. Investigate whether
plateau-edge refinement or a calibrated midpoint correction can improve
S\&C without sacrificing its safe-window probability; validate any
calibration on independent scenarios.

\item \textbf{Acquisition and information fairness.}
Include wider search intervals, noise-only frames and explicit packet
detection thresholds to quantify missed detections and false alarms.
Compare estimators given the same knowledge of the training waveform,
and examine alternative training lengths and sequence designs.

\item \textbf{Statistics and complexity.}
Use frame-level paired confidence intervals and additional independent
trials to characterize rare errors and uncertainty across channel
realizations. Evaluate recursive correlation updates and FFT-based
matched filtering, with measured execution time and memory use on a
specified platform. Such experiments would complement the approximate
operation counts reported in this thesis.
\end{enumerate}

These proposals have not been implemented.

\section{Concluding Remarks}

The central conclusion is that exact boundary acquisition, structurally
safe FFT-window placement and reliable bit detection are related but
distinct objectives. Under the tested conditions, the selected CAZAC
design achieves the strongest timing accuracy, while the S\&C results
show the importance of residual phase compensation. Timing statistics and
BER should therefore be evaluated together, under explicit assumptions
about the channel, prior information and receiver complexity.
